\section{A precise route to the first nonabelian Euler factors}
\label{sec:euler}

\subsection{The source deliberately uses some weaker interfaces}

The source's guide warns that several internal theories are developed
only in the strength needed by the proof. This matters when extracting
a theorem for another application. In particular,
\texttt{GaloisRep.Attached} gives coefficientwise \(p\)-adic limits along
a specially defined filter of Frobenius data. This is strong enough
for uniqueness of the attached trace. It is not an ordinary
quantitative density estimate for all primes.
At a split place carrying the specified Hecke polynomial, and away
from the coefficient prime, \texttt{IsCompatibleAt} says that this
polynomial annihilates each geometric Frobenius operator. Above the
coefficient prime, the interface instead uses a crystalline period
condition. Neither definition alone identifies the full characteristic
polynomial with all its multiplicities.\cite{attached}

An annihilating polynomial is often sufficient inside the source's
argument. A transfer asserting an exact Euler product must justify the
stronger identification when it is needed. In rank two this costs
much less than a general multiplicity theorem.

\begin{lemma}[Rank-two determinant upgrade]
\label{lem:ranktwo}
Let \(K\) be a field, \(A\in\GL_2(K)\), and
\(Q(X)=X^2-tX+d\). If
\[
 Q(A)=0,\qquad \det A=d,
\]
then \(t=\tr A\), and \(Q\) is the characteristic polynomial of \(A\).
The assertion holds in every characteristic.
\end{lemma}
\begin{proof}
Cayley--Hamilton gives
\(A^2-(\tr A)A+(\det A)I=0\).
Subtract \(Q(A)=0\) and use \(\det A=d\):
\((\tr A-t)A=0\). Multiplication by \(A^{-1}\) proves the claim.
\end{proof}

\subsection{The conductor-23 target}

The polynomial \(X^3-X-1\) has no rational root and discriminant
\(-23\). Its splitting field therefore has Galois group \(S_3\).
The two-dimensional standard representation is the complementary
summand to the constants in the permutation representation on the three
roots. Its determinant is the sign character, which here is the
quadratic character \(\chi_{-23}\).
At an unramified prime, let \(a_p\) denote its trace.

\begin{center}
\begin{tabular}{ccc}
\toprule
Factor degrees of the cubic modulo \(p\) & \(a_p\) & \(\chi_{-23}(p)\)\\
\midrule
\(1+1+1\) & \(2\) & \(1\)\\
\(1+2\) & \(0\) & \(-1\)\\
\(3\) & \(-1\) & \(1\)\\
\bottomrule
\end{tabular}
\end{center}

These values follow just by counting fixed roots and subtracting one.
The corresponding local Euler denominator is
\[
 1-a_p p^{-s}+\chi_{-23}(p)p^{-2s}.
\]
In a proposed theta-lift realization, suppose the polynomial
\(X^2-a_pX+\chi_{-23}(p)\) is known to annihilate the intended
Frobenius operator and its determinant has the stated sign value.
Lemma~\ref{lem:ranktwo} gives the exact local factor.
Arithmetic and geometric Frobenius have the same trace and determinant
in this \(S_3\) representation, so the convention does not change the factor.
Thus this \emph{particular interface obligation} reduces to determinant
compatibility. It does not require a new general-rank theorem.

There are two important boundaries. First, the annihilator must be the
one attached to the intended representation in the new construction;
the lemma does not identify unrelated Hecke operators. Second, the
Artin representation and its determinant are known abstractly, but the
new arithmetic lattice whose divisibility detects those \(a_p\) is still
missing. Establishing exact Euler factors does not prove a zero-free
half-plane.

\subsection{Why the rank-two simplification cannot be extrapolated}

Over \(\Q\), the two matrices
\[
 A=I_3,\qquad
 B=\begin{pmatrix}0&0&1\\1&0&0\\0&1&0\end{pmatrix}
\]
are orthogonal, semisimple, and of finite order, with determinant one.
Both are annihilated by the squarefree polynomial \(X^3-1\), but
\[
 \det(XI-A)=(X-1)^3,\qquad
 \det(XI-B)=X^3-1.
\]
Their traces are \(3\) and \(0\). Even the correct rank, determinant,
finite image, and a squarefree annihilator do not control multiplicities
in rank three. This is a counterexample to an abstract inference,
not to any theorem asserted by the upstream source.

There is a sufficient replacement in any rank \(d\): if \(A\) has a
cyclic vector \(v\), so that \(v,Av,\ldots,A^{d-1}v\) form a basis,
its minimal polynomial has degree \(d\). Every monic degree-\(d\)
annihilator is then its characteristic polynomial. A nonzero Krylov
determinant is a finite certificate. However, a scalar Frobenius
operator in rank greater than one has no cyclic vector, so this
condition cannot be assumed uniformly over all Frobenius classes.

\subsection{The remaining nonabelian work is arithmetic}

For the conductor-23 target, a full extension still requires:
\begin{enumerate}
\item A family of integral objects with a local divisibility formula
that detects the \(S_3\) trace or the appropriate nonnegative
Frobenius-weight combination.
\item Archimedean norm estimates and a second moment that preserve the
near-full spanning proportion and keep completion costs small.
\item Determinant identities and a globally compatible family of twists,
with all local and global denominator losses inside the allowed budget.
\item A final one-sided estimate, uniform over the required scales,
that reaches the square-root threshold for the correct prime remainder.
\end{enumerate}
The trace-depth theorem improves a specific input to item 3.
The rank-two lemma simplifies a local identification needed before
item 4. Neither supplies items 1 or 2. This is the concrete reason the
present work stops short of nonabelian RH.
